\section*{第 \ref{ch:b1:proteins} 章——氨基酸与蛋白质}

\begin{solution}{exo:b1:proteins:1}
$\mathrm{^+H_3N{-}CH(R){-}COO^-}$：\omterm{def:b1:proteins:aminoacid}{两性离子}。Leu 为非极性
\omterm{def:b1:lipids:triglyceride}{脂肪}族；Ser 为极性不带电；Lys 带正电；Glu 带负
电；Phe 为芳香族。
\end{solution}

\begin{solution}{exo:b1:proteins:2}
一级结构：\omterm{def:b1:proteins:aminoacid}{氨基酸}序列，由\omterm{def:b1:proteins:peptide}{肽键}连接。\omterm{def:b1:proteins:secondary}{二级结构}：局部主链
折叠（螺旋、折叠片），靠主链\omterm{def:b1:water-small-molecules:water}{氢键}维持。\omterm{def:b1:proteins:tertiary}{三级结构}：整条链的
折叠方式；\omterm{prop:b1:water-small-molecules:properties}{疏水效应}、\omterm{def:b1:water-small-molecules:water}{氢键}、盐桥、范德
瓦耳斯接触、\omterm{def:b1:proteins:tertiary}{二硫键}。\omterm{def:b1:proteins:tertiary}{四级结构}：亚基的组装，
依靠相同的非共价作用力。
\end{solution}

\begin{solution}{exo:b1:proteins:3}
$450\times 110 = \qty{50}{kDa}$。螺旋：$450\times 0.15 = \qty{68}{nm}$。
折叠链：$450\times 0.35 = \qty{158}{nm}$。
\end{solution}

\begin{solution}{exo:b1:proteins:4}
约 $0.98$ 与 $0.76$：负载的五分之一（\qty{22}{\%}）
被输送。
\end{solution}

\begin{solution}{exo:b1:proteins:5}
基团：N 端氨基（9.5）、Lys \omterm{def:b1:proteins:aminoacid}{侧链}（10.5）、Asp \omterm{def:b1:proteins:aminoacid}{侧链}
（3.9）、C 端羧基（2）。\omterm{def:b1:water-small-molecules:ph}{pH} 7：$+1 + 1 - 1 - 1 = 0$。\omterm{def:b1:water-small-molecules:ph}{pH} 2：
$+1 + 1 + 0 - 0.5 = +1.5$（末端羧基半电离）。\omterm{def:b1:water-small-molecules:ph}{pH}
12：$0 + 0 - 1 - 1 = -2$（Lys 大部分已去质子化）。净电荷为零
介于 Asp 的 p$K_a$（3.9）与氨基（9.5）之间，
此区间内电荷始终为 $0$：p$I \approx 6.7$（3.9 与 9.5 的平均值）。
\end{solution}

\begin{solution}{exo:b1:proteins:6}
脯氨酸的\omterm{def:b1:proteins:aminoacid}{侧链}与自身的氮原子相连，该氮原子
不再带有可用于形成螺旋\omterm{def:b1:water-small-molecules:water}{氢键}的氢，且 $\phi$ 被
锁定在约 $-60^\circ$，这一取值仅在
螺旋的第一圈可以容忍。甘氨酸\omterm{def:b1:proteins:aminoacid}{侧链}仅为一个氢，
不存在空间冲突，可采取其他残基都禁止的 $\phi$、$\psi$ 值，
包括紧密转角中的那些取值。
\end{solution}

\begin{solution}{exo:b1:proteins:7}
八个半胱氨酸的可能配对数为 $7\times 5\times 3\times 1 = 105$；随机
重新氧化得到天然配对的概率约为 \qty{1}{\%}。此外，
几乎完全恢复活性表明，决定配对的是肽链而非偶然，
即序列编码了折叠方式，\omterm{def:b1:proteins:tertiary}{二硫键}
只是将其锁定。
\end{solution}

\begin{solution}{exo:b1:proteins:8}
$P_{50}^{2.8} = 33.4$。当 2 时：$2^{2.8} = 6.96$，$Y = 0.17$；当 3.5 时：
$0.50$；当 8 时：$8^{2.8} = 337$，$Y = 0.91$。取 $n = 1$ 时：$Y =
P/(3.5 + P)$：$0.36$、$0.50$、$0.70$。从 8 到 2 的输送量：
有\omterm{def:b1:proteins:allostery}{协同性}时为 $0.74$，无\omterm{def:b1:proteins:allostery}{协同性}时为 $0.34$：具有\omterm{def:b1:proteins:allostery}{协同性}的\omterm{def:b1:proteins:peptide}{蛋白质}
输送的量是后者的两倍。
\end{solution}

\begin{solution}{exo:b1:proteins:9}
由四个相同的 \qty{50}{kDa} 亚基组成的四聚体。进行 SDS 凝胶电泳时
不加还原剂：若条带出现在 100、150 或 \qty{200}{kDa} 处，
说明亚基之间由\omterm{def:b1:proteins:tertiary}{二硫键}相连；若仅余 \qty{50}{kDa} 条带，
则亚基以非共价方式结合。
\end{solution}

\begin{solution}{exo:b1:proteins:10}
谷氨酸带电并在表面水合；缬氨酸
疏水，形成一块被水排斥的黏性区域。在
脱氧的 T 态下，邻近分子上暴露出一个互补的疏水口袋，
因此分子首尾相连，形成纤维，
使红\omterm{def:b1:cell-unit-of-life:cell}{细胞}变形。在肺部，R 态遮蔽了该口袋，
纤维溶解；在\omterm{def:b1:body-plans-tissues:tissue}{组织}中，血红蛋白处于
脱氧状态时纤维便会形成——\omterm{def:b1:cell-unit-of-life:cell}{细胞}因此呈镰刀状，堵塞那里的小
血管。
\end{solution}

\begin{solution}{exo:b1:proteins:11}
在 \qty{37}{\celsius} 时，$RT = \qty{2.58}{kJ/mol}$：$K = e^{-40/2.58} =
e^{-15.5} = \num{1.8e-7}$：每五百万个分子中只有一个未折叠。在
\qty{45}{\celsius}、\qty{10}{kJ/mol} 条件下，$RT = 2.64$：$K = e^{-3.8} =
0.022$：\qty{2}{\%} 未折叠，且迅速上升。临界稳定性
使\omterm{def:b1:proteins:peptide}{蛋白质}在工作时能够改变形状（别构、催化），
可为跨膜运输而展开，并在受损或不再需要时被降解、
迅速替换。
\end{solution}

\begin{solution}{exo:b1:proteins:12}
安芬森证明，序列足以决定折叠方式，
\omterm{def:b1:proteins:allostery}{别构作用}表明，形状——及其变化——即是功能。
但\omterm{def:b1:cell-unit-of-life:cell}{细胞}在每一步都会介入：\omterm{prop:b1:proteins:denaturation}{分子伴侣}挽救那些
会在折叠前聚集的肽链，因为拥挤的\omterm{def:b1:eukaryotic-cell:organelle}{细胞溶胶}并非
稀释的试管；单个替换（镰状\omterm{def:b1:cell-unit-of-life:cell}{细胞}病）得到一个
折叠正确但表面有误的\omterm{def:b1:proteins:peptide}{蛋白质}，因此“形状”必须
包括表面化学性质，而不只是主链；并且许多
\omterm{def:b1:proteins:peptide}{蛋白质}在折叠后还需添加修饰、伴侣或配体，
才能发挥作用。这句话陈述了原理，而\omterm{def:b1:cell-unit-of-life:cell}{细胞}提供
所需的条件。
\end{solution}

\begin{solution}{pb:b1:proteins:1}
\textbf{1.} $150\times 1.34 = \qty{201}{mL}$：每升血液中 $\mathrm{O_2}$ 的量。
\textbf{2.} $e^{3.51} = 33.4$；$e^{4.67} = 106.7$；$e^{7.25} = 1408$。
\textbf{3.} 肺：$1408/(33.4 + 1408) = 0.977$。\omterm{def:b1:body-plans-tissues:tissue}{组织}：
$106.7/(33.4 + 106.7) = 0.762$。
\textbf{4.} $0.977 - 0.762 = 0.215$：$0.215\times 201 = \qty{43}{mL}$
每升。
\textbf{5.} $250/43 = \qty{5.8}{L/min}$——接近静息时的
心输出量。
\textbf{6.} 肌红蛋白：肺 $13.3/13.67 = 0.973$；\omterm{def:b1:body-plans-tissues:tissue}{组织}
$5.3/5.67 = 0.935$；输送 $0.038\times 201 = \qty{7.7}{mL}$ 每
升，少了六倍：双曲线型载体装载良好，却
无法在\omterm{def:b1:body-plans-tissues:tissue}{组织}压强下释放。
\textbf{7.} \omterm{def:b1:proteins:allostery}{协同性}使曲线的陡峭部分位于
两个工作压强之间，因此压强的微小下降即可释放
很大比例的负载。
\textbf{8.} $e^{2.78} = 16.1$；$Y = 16.1/(33.4 + 16.1) = 0.325$。
\textbf{9.} $e^{4.21} = 67.4$；$Y = 16.1/(67.4 + 16.1) = 0.193$。
\textbf{10.} 无玻尔效应时：$(0.977 - 0.325)\times 201 = \qty{131}{mL}$；
有玻尔效应时：$(0.977 - 0.193)\times 201 = \qty{158}{mL}$，每升；增益
\qty{27}{mL}，\qty{20}{\%}。
\textbf{11.} $2.7/3.07 = 0.88$，已饱和；在 \qty{0.5}{kPa} 时，
$0.5/0.87 = 0.57$：收缩时，它将储备的三分之一释放到
\omterm{def:b1:eukaryotic-cell:mitochondrion}{线粒体}中，并在两次收缩之间重新装载。
\textbf{12.} 质子结合于远离血红素的位点（组氨酸、N 端），
结合后稳定 T 态，降低
血红素处的亲和力：一个位点的配体改变了
另一位点的亲和力——即\omterm{def:b1:proteins:allostery}{别构作用}。
\textbf{13.} $e^{5.45} = 233$；肺 $233/(33.4 + 233) = 0.875$；
\omterm{def:b1:body-plans-tissues:tissue}{组织} $0.762$；输送 $0.113\times 201 = \qty{23}{mL}$：约为
海平面时的一半。
\textbf{14.} $e^{3.88} = 48.4$：肺 $233/(48.4 + 233) = 0.828$；
\omterm{def:b1:body-plans-tissues:tissue}{组织} $106.7/(48.4 + 106.7) = 0.688$；输送 $0.140\times 201 =
\qty{28}{mL}$。
\textbf{15.} 装载量略有下降（由 0.875 到 0.828），而
释放量下降更多（由 0.762 到 0.688），差值因此增大：在
\qty{7}{kPa} 时，肺仍处于曲线的平坦顶部，而
\omterm{def:b1:body-plans-tissues:tissue}{组织}则位于其陡峭部分。
\textbf{16.} 容量 $180\times 1.34 = \qty{241}{mL}$；$0.140\times
241 = \qty{34}{mL}$，每升——恢复了海平面时的四分之三。
\textbf{17.} $6^{2.8} = e^{2.8\times 1.79} = e^{5.02} = 151$：肺
$233/384 = 0.607$，\omterm{def:b1:body-plans-tissues:tissue}{组织} $106.7/258 = 0.414$：输送 $0.193$，
纸面上更好——但动脉血的饱和度将仅为
\qty{61}{\%}，此后任何进一步的肺压下降，或
对更高\omterm{def:b1:body-plans-tissues:tissue}{组织}压强的任何需求，都会使大脑供氧不足；
实际的调节止于 \qty{4}{kPa} 附近。
\textbf{18.} $48.4/(33.4 + 48.4) = 0.59$。
\textbf{19.} $e^{2.68} = 14.6$；$48.4/(14.6 + 48.4) = 0.77$。
\textbf{20.} 在相同压强下，胎儿血红蛋白结合得更多：随着
母体血液饱和度降至 \qty{59}{\%}，胎儿血液
升至 \qty{77}{\%}，氧气顺着亲和力差异在胎盘膜两侧维持的
压强梯度流动。
\textbf{21.} BPG 结合 T 态并降低亲和力；$\gamma$
链与之结合较弱，因此胎儿血红蛋白更接近其
固有的高亲和力状态，$P_{50}$ 较低。
\textbf{22.} 肺：$1408/(14.6 + 1408) = 0.990$；\omterm{def:b1:body-plans-tissues:tissue}{组织} $106.7/(14.6
+ 106.7) = 0.880$；输送 $0.110\times 201 = \qty{22}{mL}$，为
成体的一半：从母体获取氧气时有利，但
肺开始工作后输送氧气则不利，因此需要转换。
\textbf{23.} $5^{2.8} = e^{2.8\times 1.609} = e^{4.51} = 90.6$：肺
$1408/1499 = 0.94$，\omterm{def:b1:body-plans-tissues:tissue}{组织} $106.7/197 = 0.54$：输送 $0.40$，几乎
为正常值的两倍——此人嘴唇轻度发绀，但
耐受良好；代价是在高海拔时储备较小。
\textbf{24.} $Y = P/(3.5 + P)$：肺 $0.79$，\omterm{def:b1:body-plans-tissues:tissue}{组织} $0.60$：输送
$0.19\times 201 = \qty{38}{mL}$——而在冲刺时的肌肉中，在
\qty{2.7}{kPa} 时，$0.44$：输送 $0.35\times 201 = \qty{70}{mL}$，与
\qty{158}{mL} 相比。无\omterm{def:b1:proteins:allostery}{协同性}时，载体的装载和
释放都很差。
\textbf{25.} 静息时每升输送 \qty{43}{mL} 的 $\mathrm{O_2}$，
冲刺肌肉中每升输送 \qty{158}{mL}；差异来自
\omterm{def:b1:proteins:allostery}{协同性}（S 形曲线）以及玻尔效应（
酸性条件下 $P_{50}$ 的位移）。
\end{solution}
